% !TEX root = ../frode-bibelen.tex
% ----------------------------------------------------------------------------
%  Frode-bibelen -- 5 Funksjoner og parametre
%  Inkluderes fra frode-bibelen.tex (kompiler alltid hovedfilen).
% ----------------------------------------------------------------------------

% ====================================================== 5 FUNKSJONER ========
\kapittel{Funksjoner og parametre}

\textbf{Prototypen} (deklarasjonen) står før \cc{main}, \textbf{definisjonen} etter. Typene må stemme overens.

\begin{center}
\begin{tabular}{@{}L{4.8cm}L{4.6cm}L{6.2cm}@{}}
\toprule
\textbf{Parameter skrevet som \dots} & \textbf{Funksjonen får \dots} & \textbf{Endring inni funksjonen \dots} \\
\midrule
\cc{int n}, \cc{char c}, \cc{float f} & en \textbf{kopi} & påvirker \emph{ikke} kalleren \\
\cc{const int n} & en kopi som ikke kan endres & -- \\
\cc{int* n} (kall med \cc{\&x}) & adressen (\emph{referanse}) & endrer \textbf{originalen} via \cc{*n = \dots} \\
\cc{int a[]} / \cc{char t[]} / \cc{char* t} & adressen til arrayen & endrer \textbf{originalen} (arrayer kopieres aldri) \\
\cc{const char* t} & adressen, skrivebeskyttet & ikke lov -- brukes for tekster som bare leses \\
\cc{struct Kunde k} & en \textbf{kopi} av hele structen & påvirker \emph{ikke} originalen \\
\cc{struct Kunde* k} & adressen til structen & endrer \textbf{originalen} (\cc{k->saldo = 0;}) \\
\cc{const struct Kunde* k} & adressen, skrivebeskyttet & ikke lov -- for skriv-/sjekk-funksjoner \\
\bottomrule
\end{tabular}
\end{center}

\subsection{Returverdier}
\begin{kode}
bool kundeErType(const struct Kunde* kunde, const enum Type type)  {
  return (kunde->type == type);              //  Returnerer sant/usant.
}

int finnKunde(const char* navn)  {           //  Returnerer INDEKS eller -1:
  for (int i = 0;  i < gAntallKunder;  i++)
      if (!strcmp(gKunder[i]->navn, navn))  return i;
  return -1;
}

float kundeHentSaldo(const struct Kunde* kunde)  {  return kunde->saldo;  }
\end{kode}

\subsection{Flere «returverdier» -- referanseoverføring med pekere}
En funksjon kan bare \cc{return}-e én verdi. Trenger du flere, send \textbf{adresser} inn (S24 brukte dette):
\begin{kode}
void kundeHentTall(const struct Kunde* kunde, int* tlf, float* saldo)  {
  *tlf   = kunde->tlf;                       //  Skriver til kallerens variable.
  *saldo = kunde->saldo;
}
                                             //  Kall:
int t;   float s;
kundeHentTall(gKunder[i], &t, &s);           //  NB:  & foran!
\end{kode}

\subsection{Bytte om to}
\begin{kode}
void bytt(int* a, int* b)  {  int temp = *a;   *a = *b;   *b = temp;  }
bytt(&x, &y);

struct Kunde* hjelp = gKunder[i];            //  Bytte to PEKERE i arrayen
gKunder[i] = gKunder[j];   gKunder[j] = hjelp;   //  (structene flyttes ikke).

void byttStruct(struct Kunde* s1, struct Kunde* s2)  {   //  Bytte INNHOLD:
  struct Kunde s = *s1;   *s1 = *s2;   *s2 = s;
}
\end{kode}

\begin{felle}[{Parameter-feller}]
\begin{itemize}
  \item Glemt \cc{\&} i kallet: \cc{kundeHentTall(k, t, s)} gir kompileringsfeil / krasj.
  \item Endrer du en \cc{struct}-\emph{kopi} (\cc{struct Kunde k}) inni funksjonen, ser ikke kalleren endringen. Klassisk i oppgave 1b!
  \item Arrayparametre har ingen størrelse -- send med antallet: \cc{void skriv(int a[], int n)}.
\end{itemize}
\end{felle}
